% =====================================================================
%  IEEE Journal of Photovoltaics -- main manuscript
%  Built on the official IEEEtran class shipped in
%  IEEE-Transactions-LaTeX2e-templates-and-instructions.
%
%  BUILD:  ./build.sh          (or: pdflatex, bibtex, pdflatex, pdflatex)
%
%  EDITORIAL NOTES marked \edit{...} are typeset in red and MUST be removed
%  before submission. Set \editsfalse below to hide them all at once.
% =====================================================================
\documentclass[lettersize,journal]{IEEEtran}

\usepackage{amsmath,amsfonts,amssymb}
\usepackage{array}
\usepackage{textcomp}
\usepackage{stfloats}
\usepackage{url}
% let long repository URLs break rather than run into the margin
\def\UrlBreaks{\do\/\do-\do.\do:\do_}
\usepackage{graphicx}
\usepackage{cite}
\usepackage{xcolor}
\usepackage{booktabs}

% ---- editorial flag machinery -------------------------------------------
\newif\ifedits
\editstrue          % <-- set to \editsfalse to strip every editorial note
\newcommand{\edit}[1]{\ifedits{\color{red}\textbf{[#1]}}\fi}

% ---- convenient chemistry shorthands ------------------------------------
\newcommand{\FASnI}{FASnI\textsubscript{3}}
\newcommand{\CABI}{Cu\textsubscript{2}AgBiI\textsubscript{6}}
\newcommand{\BZS}{BaZrS\textsubscript{3}}
\newcommand{\CABB}{Cs\textsubscript{2}AgBiBr\textsubscript{6}}
\newcommand{\TiOx}{TiO\textsubscript{2}}
\newcommand{\percm}{\,mA\,cm\textsuperscript{$-$2}}

\begin{document}

% NOTE ON THE TITLE. The subtitle currently advertises the 3-D full-wave
% method. Every result in this paper is planar, and a planar stack at normal
% incidence with periodic lateral boundaries is equivalent to a 1-D coherent
% calculation, so a referee may reasonably ask what the 3-D bought. Consider
% dropping the subtitle, or replacing it with
% ``A Controlled Opto-Electro-Thermal Comparison with Coherent Optics''.
% Previous title kept for reference:
%   ...: A Full-Wave 3-D Opto-Electro-Thermal Study
% Subtitle dropped: every result is planar, so the 3-D framework is presented
% as machinery VALIDATED against the transfer-matrix limit (Secs. II-B, II-E)
% rather than as the contribution.
\title{Interface Band Offsets, Not Optical Collection, Limit Four Chemically
Distinct Lead-Free Absorbers}

\author{Md.~Mohiuddin and Alamgir~Kabir%
\thanks{The authors are with the Department of Physics, University of Dhaka,
Dhaka 1000, Bangladesh (e-mail: alamgir.kabir@du.ac.bd).}%
\thanks{Manuscript received \edit{DATE}; revised \edit{DATE}.}}

\markboth{IEEE Journal of Photovoltaics,~Vol.~XX, No.~X, Month~2026}%
{Mohiuddin \MakeLowercase{\textit{et al.}}: Interface Band Offsets Limit Four Lead-Free Absorbers}

\maketitle

\begin{abstract}
Lead-free absorbers are the leading route to non-toxic halide-perovskite
photovoltaics, but their device-simulation literature is dominated by
one-dimensional, single-material optimizations that prescribe the optical
generation profile, fix poorly known parameters, and rarely disclose parameter
provenance, yielding projected efficiencies difficult to reconcile with
experiment. We report a unified opto-electro-thermal finite-element study of four chemically distinct lead-free absorbers (\BZS,
\FASnI, \CABB, \CABI) in a single fixed architecture, coupling full-wave
electromagnetic optics to drift--diffusion transport, with every fixed
parameter literature-sourced and provenance-graded and every uncertain
parameter swept. Solving Maxwell's equations on the real multilayer stack
rather than assuming Beer--Lambert absorption changes the optical current
ceilings by $-16\,\%$ to $+18\,\%$ relative to prescribed-profile estimates,
and in one case the previously reported device photocurrent exceeded the
corrected ceiling, a physically impossible result that the full-wave treatment
resolves. The electromagnetic solver is verified against an independent
transfer-matrix calculation to better than $0.1\,\%$ absorptance at every
wavelength, a check that also exposed and corrected a back-contact truncation
artifact in one absorber. Under defect-free interfaces no device is limited by optical
collection: the best collects $99.6\,\%$ of its optical ceiling. Each is
instead limited in voltage and fill factor by a band-offset cliff at a
transport interface. The four materials partition into two mechanistic
classes, electron-transport-layer conduction-band cliffs (\FASnI, \CABI) and
hole-transport-layer valence-band cliffs (\BZS, \CABB). Closing the thermal
loop for \CABB{} gives a benign relative temperature coefficient of
$\approx -0.12\,\%\,\mathrm{K}^{-1}$. Parasitic absorption in the transparent
conductor is a first-order loss throughout, exceeding the absorber's own
photocurrent in the most weakly absorbing case. Transport-layer band
engineering, not further absorber optimization, is the operative design
problem for this class.
\end{abstract}

\begin{IEEEkeywords}
Lead-free perovskite, photovoltaic cells, finite-element methods, numerical
simulation, band offset, opto-electro-thermal modeling.
\end{IEEEkeywords}

\section{Introduction}
\IEEEPARstart{E}{very} halide perovskite solar cell that has crossed
$26\,\%$ certified power-conversion efficiency (PCE) is built on an absorber
whose defining component, divalent lead, is water-soluble and toxic, an
unresolved liability for the terawatt-scale, environmentally exposed
deployment the technology is meant to serve~\cite{nrel_chart}. Regulatory
pressure and life-cycle concerns have made the replacement of lead the central
materials problem of the field, and a chemically diverse set of lead-free
absorbers has emerged in response~\cite{prakash2026}. These candidates are not
variations on a theme: they span fundamentally different bonding,
dimensionality, defect chemistry, and, crucially for device design, band-edge
energetics. That diversity is an opportunity, but it also means lessons from
optimizing one lead-free absorber rarely transfer to another.

Four families frame the landscape. \emph{Tin halide perovskites} such as
\FASnI{} are the closest structural analogue to the lead compounds and the
most advanced in single-junction devices~\cite{fasni3_14pct}, reaching
$14.51\,\%$ certified on $1\,\mathrm{cm}^2$ fullerene-free
architectures~\cite{li2026} and $16.65\,\%$ through 2D/3D heterostructure
engineering~\cite{he2025}, though Sn$^{2+}\!\rightarrow$Sn$^{4+}$ oxidation
remains their governing stability problem~\cite{yang2017}.
\emph{Chalcogenide perovskites} such as \BZS{} replace halide with sulfide,
offering exceptional thermal and moisture stability~\cite{kumar2021}; their
photovoltaic literature, however, remains overwhelmingly computational, with
SCAPS-1D optimizations projecting $22$--$34\,\%$
efficiencies~\cite{sowayan2025,guo2025} that far outrun nascent experimental
demonstrations. \emph{Halide double perovskites} such as \CABB{} restore full
inorganic stability, but wide indirect gaps and sluggish transport cap
experimental single-junction PCE near
$6\,\%$~\cite{zhang2022,li2021review,moed2026} (a record that itself required
gap-narrowing hydrogenation; unmodified devices sit nearer $3$--$4\,\%$
\cite{zhang2022}). \emph{Quaternary
bismuth/silver iodides} such as \CABI{} are strongly absorbing, air-stable,
wide-gap semiconductors attractive for tandem top cells and indoor
harvesting~\cite{sansom2021,moon2025}, yet their experimental efficiencies
remain near $2\,\%$~\cite{islam2023,pai2022}. Across all four the
pattern repeats: promising absorber physics, modest real devices, and a
simulation literature whose optimism is hard to reconcile with experiment.

That reconciliation problem is largely methodological. The device-simulation
studies populating this literature are predominantly one-dimensional,
single-material optimizations that (i) prescribe the optical generation
profile from an assumed absorption coefficient rather than solving Maxwell's
equations for the real multilayer stack, (ii) fix critical and poorly known
parameters at single borrowed values and then optimize around them,
(iii) treat the cell as isothermal at $300\,\mathrm{K}$, and (iv) rarely trace
the origin of the inputs that determine the answer. Under these conditions it
is unsurprising that an absorber capped near $6\,\%$ in the laboratory can be
projected above $25\,\%$ on paper.

This work addresses those gaps with a unified opto-electro-thermal
finite-element study of the four absorbers in a fixed architecture. Full-wave electromagnetic optics feed a finite-volume
drift--diffusion transport model; every fixed material parameter is
literature-sourced with an explicit provenance status, and each genuinely
uncertain quantity is treated as a sensitivity sweep. Two results follow.
First, replacing prescribed generation profiles with full-wave optics shifts
the optical ceilings by $-16\,\%$ to $+18\,\%$, and in one case removes a
thermodynamic inconsistency in previously reported device numbers. Second,
under defect-free interfaces no device is optically limited; each is limited
by a band-offset cliff at a transport interface, and the four partition into
two mechanistic classes according to which contact carries that cliff.

\section{Computational Method}
All devices were simulated in COMSOL Multiphysics (v6.2 locally, v6.3 on the
compute server) using a single self-consistent workflow coupling full-wave
electromagnetic optics to finite-volume drift--diffusion transport, with an
additional steady-state heat-transfer problem for \CABB.

The governing principle for parameterization was to admit no unjustified point
assumptions. Every fixed material parameter is literature-sourced and carries
an explicit provenance status: \textbf{[EXP]} experimental, \textbf{[THEO]}
computed with the level stated, \textbf{[SIM-LIT]} from published device
simulations, \textbf{[DERIVED]} computed here from literature inputs through a
stated parameter-free relation. Quantities that are genuinely unknown or that
conflict across the literature (electron affinity $\chi$, carrier lifetime
$\tau$, relative permittivity $\varepsilon_r$, mobility $\mu$, doping) are
\textbf{[SWEPT]} sensitivity axes. Full status-coded tables are in the
Supplementary Material.

\subsection{Device architecture}
A fixed stack is used throughout: glass / FTO ($600\,\mathrm{nm}$) / \TiOx{}
ETL ($50\,\mathrm{nm}$) / absorber ($300\,\mathrm{nm}$) / Spiro-OMeTAD HTL
($200\,\mathrm{nm}$) / Au. The absorber is the only layer that changes between
devices, so all comparisons are on identical optical, electrical and numerical
footing. The optical unit cell is $350\times350\,\mathrm{nm}^2$ with Floquet
periodic boundaries on all four lateral faces, a perfectly matched layer
above, and a scattering boundary condition at the illuminated plane.

\subsection{Optical model}
The electromagnetic problem is solved in the frequency domain, total-field
formulation, for 61 wavelengths spanning $300$--$900\,\mathrm{nm}$ at
$10\,\mathrm{nm}$ resolution under AM1.5G illumination~\cite{astm_g173}.
Absorbed power density is integrated per layer, giving a closed optical
budget: absorber plus each parasitic layer plus reflectance equals incident.
Absorber optical constants are measured data throughout: tabulated
ellipsometric $\varepsilon$ for \FASnI~\cite{ghimire2017}, tabulated
ellipsometry with a photothermal-deflection sub-gap tail for
\CABI~\cite{kamppinen2026,sansom2021}, the measured ellipsometric dielectric
function for \BZS~\cite{nishigaki2020}, and the conductivity-route extinction
derived from~\cite{eddekkar2024} for \CABB. The contacts use measured
SnO$_2$:F dispersion~\cite{vonrottkay1996} and Au~\cite{mcpeak2015}; \TiOx{}
uses a Sellmeier index~\cite{devore1951}.

Two aspects require care and are the origin of the ceiling changes reported
here. First, \emph{Floquet periodic boundaries require the paired faces to
carry identical surface meshes.} A free-tetrahedral mesher does not produce
this, and the solver then interpolates across the mismatch without raising an
error; the computed absorptance wanders with every remesh. On this stack,
absorber absorptance at $700\,\mathrm{nm}$ took values $0.719$, $0.663$ and
$0.463$ on meshes differing by $2.5\,\%$ in element count, a $36\,\%$ swing.
The fix is to mesh one face of each pair with a free-triangular operation and
copy it onto the opposite face before the tetrahedral fill; the same
comparison then gives $0.783103$ versus $0.782888$ ($0.03\,\%$). Critically,
\emph{the energy-conservation residual does not detect this error}: it
compares total absorbed power against total incident flux and is blind to how
that total is partitioned between layers, which is exactly what the
photocurrent depends on.

Second, meshes are banded by wavelength ($300$--$390$, $400$--$590$,
$600$--$900\,\mathrm{nm}$) because element count scales as $\lambda^{-3}$;
element size is set per material as $\lambda/(n_\mathrm{eff}\,\mathrm{ppw})$
with $n_\mathrm{eff}=|\tilde{n}|$, subject to a skin-depth clamp and a floor of
three elements per layer thickness. Wavelengths at the band boundaries are
solved on both neighbouring meshes as a seam check.

The AM1.5G weighting integrates the irradiance over each wavelength bin rather
than sampling at bin centres, with outer bin edges clamped to the grid
endpoints.

For the laterally uniform stacks considered here at normal incidence, the
periodic three-dimensional problem is mathematically equivalent to a
one-dimensional coherent multilayer problem. The finite-element framework is
used because it extends unchanged to non-planar geometry; Section~II-E
verifies it against that analytic limit, which is also what exposed the one
place where a geometric idealization mattered.

\subsection{Electrical and thermal models}
Transport is finite-volume drift--diffusion with Fermi--Dirac statistics,
Shockley--Read--Hall recombination, and a radiative term
$R = B(np-n_i^2)$ whose coefficient $B$ is derived from the van
Roosbroeck--Shockley relation~\cite{vanroosbroeck1954} using the same optical
constants that drive the optical model, making it parameter-free rather than
fitted. Heterojunctions are treated in both quasi-Fermi-continuity and
thermionic-emission modes; the two agree to $0.2\,\%$. The broadband
generation profile $G(z)$ from the optical stage is laterally averaged and
supplied to the transport model. For \CABB{} a steady-state heat-transfer
problem is coupled to the optical and electrical stages, following coupled
opto-electro-thermal treatments of lead-based
devices~\cite{saxena2019,bahrami2024,suldozi2026} and motivated by the
reported sensitivity of self-heating to active-material
properties~\cite{kamppinen2024}, with per-layer thermal conductivities all
literature-cited~\cite{mun2007,oka2014,ho1972,wang2015,zheng2024}. The
thermal stage was closed on the optical inputs of the parent study; the
quoted temperature coefficient is a relative quantity and is robust to the
subsequent rescaling of the generation profile, and absolute operating
temperatures are therefore not quoted here.

\subsection{Numerical procedure and validation}
Every result carries three independent checks. The optical budget must close,
and the depth integral of $G(z)$ must reproduce the photocurrent obtained
independently from the spectral integral of layer absorptance; these are
different routes through the same solve. Across all four absorbers the budget
closes to $\pm0.00\,\%$, the worst per-wavelength energy residual is
$-1.43\,\%$, the two photocurrent routes agree within $0.05\,\%$, and band
seams agree to $0.020$--$0.096\,\%$. Third, an independent scalar
transfer-matrix calculation with identical optical inputs and identical
geometry reproduces the finite-element absorber absorptance to within
$7\times10^{-4}$ at every wavelength and the photocurrent to
$\leq0.02\,\%$, for all four absorbers (Fig.~S7) -- a different solver
checked against the analytic planar limit, not another route through the same
solve. This third check exposed a real artifact: the optical model truncates
the Au contact at $45\,\mathrm{nm}$ over a perfect electric conductor,
justified by Au's $12$--$17\,\mathrm{nm}$ near-infrared skin depth, but Au's
damping dips in its $450$--$570\,\mathrm{nm}$ interband window, and \CABB{}
is the one absorber transparent exactly there. The truncation inflated its
ceiling by $4.3\,\%$; all \CABB{} results below use the corrected
thick-contact limit, computed with the validated transfer-matrix method and
re-solved through transport (Supplementary Section~S3).

Device $J$--$V$ curves use continuation in applied bias. For \CABI{} a
$25\,\mathrm{mV}$ grid lands on spurious negative-differential-resistance
branches; a $1\,\mathrm{mV}$ grid removes them, and only the dense result is
quoted.

\section{Results}

\subsection{Full-wave optics changes the current ceilings}
Fig.~\ref{fig:optics} shows the optical constants and
Fig.~\ref{fig:budget} the resulting absorptance, reflectance and closed loss
budget. Table~\ref{tab:ceilings} compares the optical current ceiling, the
photocurrent at $100\,\%$ collection, against the previous prescribed-profile
estimates.

\begin{figure}[!t]
\centering
\includegraphics[width=\columnwidth]{figures/fig1_optical_inputs.pdf}
\caption{Optical constants of the four absorbers over
$300$--$900\,\mathrm{nm}$: (a)~refractive index $n$; (b)~extinction
coefficient $k$ (logarithmic).}
\label{fig:optics}
\end{figure}

\begin{table}[!t]
\caption{Optical Current Ceilings (mA\,cm$^{-2}$, AM1.5G, 100\,\% Collection)}
\label{tab:ceilings}
\centering
\begin{tabular}{lrrr}
\toprule
Absorber & Previous & This work & Change \\
\midrule
\FASnI  & 18.27 & \textbf{20.450} & $+11.9\,\%$ \\
\CABI   & 13.91 & \textbf{15.083} & $+8.4\,\%$ \\
\CABB   & 2.07  & \textbf{2.440}$^{\dagger}$  & $+17.9\,\%$ \\
\BZS    & 12.29 & \textbf{10.306} & $-16.1\,\%$ \\
\bottomrule
\end{tabular}

\vspace{2pt}
{\footnotesize $^{\dagger}$Thick-contact limit, correcting the Au-truncation
artifact of Sec.~II-E; the truncated value was 2.546.}
\end{table}

Three of the four ceilings rise, principally because the transparent conductor
is now described by measured dispersion~\cite{vonrottkay1996} rather than a
flat placeholder extinction coefficient. \BZS{} falls, for a different and
more consequential reason: its optical constants were previously a
Kramers--Kronig reconstruction of a \emph{computed} dielectric function that
was zero-padded at $690\,\mathrm{nm}$, inside its own absorption range,
whereas this work uses the measured ellipsometric dielectric
function~\cite{nishigaki2020}. The measured data has markedly weaker near-edge
absorption and an onset at $639\,\mathrm{nm}$.

The consequence is material. The previously reported \BZS{} device
photocurrent of $11.85$\percm{} \emph{exceeds} the corrected optical ceiling
of $10.306$\percm, which is impossible: the transport model was being supplied
generation the measured optics say does not occur.

\begin{figure}[!t]
\centering
\includegraphics[width=\columnwidth]{figures/fig2_absorptance_budget.pdf}
\caption{Full-wave optical response. (a)~Absorber absorptance and
(b)~stack reflectance versus wavelength. (c)~AM1.5G-weighted optical loss
budget. The budget closes to $\pm0.00\,\%$ of the incident
$33.74$\percm{} in every case.}
\label{fig:budget}
\end{figure}

The loss budget (Fig.~\ref{fig:budget}c) shows the parasitic cost of the
transparent conductor is first-order for every absorber, ranging from
$7.2$\percm{} (\FASnI) to $9.8$\percm{} (\CABB); in the weakly absorbing
double perovskite the contact absorbs four times more photocurrent than the
absorber does.

\subsection{Devices of record}
Every device was re-solved on the corrected generation profiles. Each re-run
was preceded by a control on the \emph{previous} profile that had to reproduce
the published base case before the new result was accepted; all four controls
reproduced their targets exactly. Fig.~\ref{fig:jv} shows the resulting
$J$--$V$ characteristics against their controls and
Table~\ref{tab:devices} gives the metrics.

\begin{figure}[!t]
\centering
\includegraphics[width=\columnwidth]{figures/fig3_generation.pdf}
\caption{Broadband generation profiles $G(z)$ through the $300\,\mathrm{nm}$
absorber, plotted from the hole-transport interface, on (a)~the previous
prescribed optics and (b)~the full-wave optics of this work.}
\label{fig:gen}
\end{figure}

\begin{figure}[!t]
\centering
\includegraphics[width=\columnwidth]{figures/fig4_jv.pdf}
\caption{Device $J$--$V$ characteristics. Dashed: control run on the previous
generation profile, which reproduces the published base case exactly. Solid:
this work, on the corrected full-wave profile.}
\label{fig:jv}
\end{figure}

\begin{table*}[!t]
\caption{Device Metrics Before and After the Full-Wave Optical Correction}
\label{tab:devices}
\centering
\begin{tabular}{lccccc}
\toprule
Absorber & $J_\mathrm{sc}$ (mA\,cm$^{-2}$) & $V_\mathrm{oc}$ (V) & FF & PCE (\%) & Collection \\
\midrule
\FASnI & $14.740\rightarrow\mathbf{15.811}$ & $0.987\rightarrow0.988$ & $0.517\rightarrow0.509$ & $7.52\rightarrow\mathbf{7.95}$  & $80.7\rightarrow77.3\,\%$ \\
\CABI  & $13.856\rightarrow\mathbf{15.017}$ & $1.379\rightarrow1.384$ & $0.519\rightarrow0.518$ & $9.91\rightarrow\mathbf{10.76}$ & $99.6\rightarrow99.6\,\%$ \\
\BZS   & $11.846\rightarrow\mathbf{9.883}$  & $1.341\rightarrow1.337$ & $0.340\rightarrow0.340$ & $5.40\rightarrow\mathbf{4.48}$  & $96.4\rightarrow95.9\,\%$ \\
\CABB  & $2.010\rightarrow\mathbf{2.371}$   & $1.710\rightarrow1.760$ & $0.287\rightarrow0.281$ & $0.98\rightarrow\mathbf{1.17}$  & $97.3\rightarrow97.2\,\%$ \\
\bottomrule
\end{tabular}

\vspace{2pt}
\footnotesize \CABI{} quoted on a matched $1\,\mathrm{mV}$ continuation grid.
\end{table*}

The \BZS{} inconsistency is resolved: $9.883$\percm{} now sits below its
$10.306$ ceiling. Collection efficiency ($95.9\,\%$) and fill factor ($0.340$)
are essentially unchanged, so the entire loss is optical, a corrected input
rather than a changed model.

\emph{No device is limited by optical collection.} \CABI{} collects
$99.6\,\%$ of its ceiling, \CABB{} $97.2\,\%$, \BZS{} $95.9\,\%$. Each is
instead limited in voltage and fill factor.

The one absorber whose collection \emph{changes} is instructive. \FASnI's
optical ceiling rose $11.9\,\%$ but its photocurrent rose only $7.3\,\%$,
because collection fell from $80.7\,\%$ to $77.3\,\%$: the corrected optics
front-load generation toward the illuminated face (Fig.~\ref{fig:gen}), and
\FASnI{} collects that less efficiently than the flatter previous profile.
\CABB, absorbing weakly and generating almost uniformly, passes its full
$+18\,\%$ optical gain through to the device. Whether an optical improvement
reaches the terminals therefore depends on where in the depth it deposits
carriers, a distinction invisible to a prescribed-profile model.

\subsection{Two classes of interface cliff}
The limiting mechanism in every device is a band-offset cliff at a transport
interface, and the four materials partition into two classes
(Fig.~\ref{fig:bands}). \FASnI{} and \CABI{} are limited by a
\emph{conduction-band cliff at the \TiOx{} electron contact}, $0.65\,$eV and
$0.93\,$eV respectively, which permits interfacial recombination and caps fill
factor. \BZS{} and \CABB{} are limited by a \emph{valence-band cliff at the
Spiro-OMeTAD hole contact}, $0.93\,$eV and $1.27\,$eV, which blocks forward
hole injection, produces an S-shaped $J$--$V$ characteristic and collapses
fill factor to $0.28$--$0.34$~\cite{minemoto2001}.

\begin{figure*}[!t]
\centering
\includegraphics[width=\textwidth]{figures/fig5_band_alignment.pdf}
\caption{Band alignment of the four devices against the vacuum level, with the
conduction- and valence-band offsets at the two transport interfaces. The
offset that limits each device is highlighted. The four partition into two
classes: electron-contact cliffs (a,~b) and hole-contact cliffs (c,~d).}
\label{fig:bands}
\end{figure*}

The partition is sharp. In each material the limiting offset exceeds the
opposing one by at least a factor of three, and in \BZS{} and \CABB{} by more
than eight, so the assignment does not depend on where the boundary between
``cliff'' and ``spike'' is drawn. It is also orthogonal to the chemistry: one
tin halide and one bismuth iodide fall in the first class, one chalcogenide
and one double perovskite in the second. Because electron affinity is a swept
rather than fixed parameter here, the classification was checked across the
swept range and is stable within it.

\subsection{Two natures of the open-circuit voltage}
An independent axis distinguishes how $V_\mathrm{oc}$ is set. For \FASnI{} the
radiative term dominates and $V_\mathrm{oc}$ is pinned near its radiative
limit. For the bismuth halides, bulk Shockley--Read--Hall recombination sets
it.

\BZS{} belongs to neither category, and correcting this is a substantive
change to the previously reported mechanism. We recomputed its radiative
coefficient from the van Roosbroeck--Shockley relation using the measured
optical constants and found $B_\mathrm{rad} = 5.37\times10^{-12}\,
\mathrm{cm^3\,s^{-1}}$, some $435\times$ smaller than the value in use.
Re-solving the device with the corrected coefficient moves $V_\mathrm{oc}$ by
$3\,\mathrm{mV}$. \BZS's $V_\mathrm{oc}$ is therefore \emph{not} at its
radiative bound, as previously stated; it is set by dark current through the
Spiro injection barrier. The previously used coefficient was also internally
inconsistent, the van Roosbroeck--Shockley integral having been taken from
$1.75\,$eV while $n_i^2$ was constructed from a $1.88\,$eV gap, making it
$\approx12\times$ too large even on its own optical data. Because the device
is insensitive to it, this does not affect any reported performance number.

\subsection{Numerical fidelity and disclosures}
Three caveats must be stated explicitly.

\emph{\CABI{} sub-gap absorption.} $38.2\,\%$ of this absorber's photocurrent
($5.77$\percm) originates below its $2.06\,$eV Tauc gap. This is not an
artifact of the electromagnetic solve: the optical file is constructed from
tabulated ellipsometry above $2.0\,$eV~\cite{kamppinen2026} and a
photothermal-deflection tail shape below it~\cite{sansom2021}, zeroed under
$1.40\,$eV. Consequently the ratio of photocurrent to the detailed-balance
ceiling at $2.06\,$eV reads $114\,\%$, which is meaningless; we therefore
report the above-gap/sub-gap split rather than a ceiling fraction. Within that
sub-gap current, $0.908$\percm{} ($6.0\,\%$ of the total) arises from a
segment where the absorption coefficient stops decaying and sits flat at
$\approx2.7\times10^{3}\,\mathrm{cm^{-1}}$ for $130\,\mathrm{nm}$ before
stepping to zero, a construction artifact rather than a physical Urbach tail.
An optical ceiling credits all of these carriers at $100\,\%$ collection,
which is optimistic for tail states. Re-solving the device on restricted
generation profiles bounds the assumption (Supplementary Section~S7): with
the artifact removed but the measured tail kept, PCE is $10.10\,\%$; with no
sub-gap collection at all, $6.98\,\%$. The full-collection $10.76\,\%$ of
Table~II is therefore the optimistic end of a $[6.98, 10.76]\,\%$ bracket
whose defensible central value is $10.10\,\%$.

\emph{\CABI{} open-circuit voltage is not grid-converged.} The same model on a
mixed $1$--$25\,\mathrm{mV}$ continuation grid gives $V_\mathrm{oc} =
1.512\,$V and on a uniform $1\,\mathrm{mV}$ grid $1.379\,$V, a $133\,$mV
spread. Maximum power, $V_\mathrm{mpp}$ and $J_\mathrm{mpp}$ are identical to
five significant figures ($9.9126\,\mathrm{mW\,cm^{-2}}$, $0.786\,$V,
$12.612$\percm), so PCE is unaffected. The cause is physical: the ETL cliff
makes the pre-$V_\mathrm{oc}$ tail nearly flat, so a $0.08$\percm{} numerical
difference moves the zero crossing enormously. $V_\mathrm{oc}$ is poorly
conditioned for this absorber; $P_\mathrm{max}$-derived metrics are robust.
A related gap-choice test (Supplementary Section~S2) re-solved \CABB{} at its
$2.21\,$eV optical onset and \BZS{} at its $1.94\,$eV measured onset: PCE is
invariant to four decimals in both, with the shift appearing entirely in
$V_\mathrm{oc}$ (${\approx}\Delta E_g$) and compensated by the fill factor --
the disclosed gap inconsistencies do not affect any efficiency conclusion.

\emph{Lossless transport layers.} \TiOx{}~\cite{devore1951} and Spiro-OMeTAD
are modeled with real-only refractive indices ($k=0$), so their budget entries
are exactly zero and all front parasitic absorption is attributed to the FTO.
\TiOx{} genuinely absorbs below $\approx390\,\mathrm{nm}$, so UV generation in
the absorber is slightly overestimated. The treatment is identical across all
four devices, so comparisons are unaffected.

\section{Discussion}
Three design consequences follow.

\emph{Transport-layer band engineering is the operative problem.} With
defect-free interfaces every device already collects $96$--$99.6\,\%$ of its
available photocurrent; further absorber optimization cannot address what
limits them. The cliff taxonomy identifies which contact to change for which
material: \FASnI{} and \CABI{} need a shallower-affinity electron contact,
\BZS{} and \CABB{} a deeper-HOMO hole contact. For \BZS{} neither
Spiro-OMeTAD~\cite{rahman2023} nor CuO~\cite{elnaggar2024,ungeheuer2024}
provides both a deep HOMO ($\approx5.9\,$eV) and a high conduction-band
minimum, which is why that material's HTL selection remains unresolved.

\emph{Optical improvements must be evaluated in depth, not in aggregate.} The
\FASnI/\CABB{} contrast shows that identical fractional gains in absorbed
photons produce different device gains depending on where carriers are
generated. Any light-management claim benchmarked only on absorptance may
overstate its device benefit.

\emph{Parasitic absorption in the transparent conductor is first-order.} It
costs $7.2$--$9.8$\percm{} across the four devices, and for \CABB{} the
contact absorbs four times more photocurrent than the absorber does. Any
optical improvement to these stacks must be assessed against this loss
channel, not against reflection alone.

\emph{Limitations.} Interfaces are defect-free, so reported efficiencies are
upper bounds; no literature interface-recombination data exists for these
material pairs, and junction recombination is not representable by the
solver's surface-recombination feature in this formulation (an earlier
sensitivity claim to the contrary is withdrawn and corrected in Supplementary
Section~S8). Because the limiting mechanism is carrier accumulation at a
band-offset cliff, the upper-bound label is load-bearing. Transport layers
are optically lossless. The thermal loop is closed
only for \CABB. \CABI's $V_\mathrm{oc}$ carries the convergence caveat above,
and its optical constants include a sub-gap tail whose device consequences we
do not model as distinct from band-edge carriers.

\section{Conclusion}
A unified opto-electro-thermal study of four chemically distinct lead-free
absorbers in a fixed architecture, with coherent full-wave optics verified
against the transfer-matrix limit, shows that
none is limited by optical collection under defect-free interfaces; each is
limited by a band-offset cliff at a transport interface, and the four
partition into electron-contact and hole-contact cliff classes. Solving
Maxwell's equations rather than prescribing generation changes the optical
ceilings by $-16\,\%$ to $+18\,\%$ and removes a case in which a reported
device photocurrent exceeded its own optical ceiling. Parasitic absorption in
the transparent conductor is a first-order loss in every device.
Transport-layer band engineering, not further absorber optimization, is the
operative next step for this class of materials.

\section*{Data Availability}
All simulation inputs, build scripts, result files and figure-generation code
are available at
\url{https://github.com/mohiuddinsciencelover-glitch/leadfree-pv-3d-simulation},
together with \texttt{REPRODUCE.md}, which documents the full pipeline end to
end.

\section*{Acknowledgment}
The authors thank the Department of Physics, University of Dhaka, for
computational resources.

% IEEEtran.bst gives correct IEEE reference formatting and ships with
% texlive-publishers. ieeetr.bst is a close fallback so this always builds.
\IfFileExists{IEEEtran.bst}{\bibliographystyle{IEEEtran}}{\bibliographystyle{ieeetr}}
\bibliography{references}

\vfill

% ---- author biographies -------------------------------------------------
% Replace the placeholder photographs with real files. IEEE specification:
% 1 in x 1.25 in, 300 dpi minimum, TIFF or high-quality JPEG, head and
% shoulders. Put them beside this .tex as mohiuddin.jpg / kabir.jpg and
% switch the \begin{IEEEbiographynophoto} lines below to \begin{IEEEbiography}
% with the \includegraphics argument shown in the comment.

% \begin{IEEEbiography}[{\includegraphics[width=1in,height=1.25in,clip,keepaspectratio]{mohiuddin}}]{Md. Mohiuddin}
\begin{IEEEbiographynophoto}{Md. Mohiuddin}
\edit{PHOTO REQUIRED: mohiuddin.jpg, 1$\times$1.25 in, 300 dpi}
\edit{received the B.Sc. and M.Sc. degrees in physics from the University of
Dhaka, Dhaka, Bangladesh, in YYYY and YYYY, respectively. He is currently
\ldots{} with the Department of Physics, University of Dhaka. His research
interests include computational photovoltaics, finite-element device
simulation, and lead-free perovskite absorbers.}
\end{IEEEbiographynophoto}

% \begin{IEEEbiography}[{\includegraphics[width=1in,height=1.25in,clip,keepaspectratio]{kabir}}]{Alamgir Kabir}
\begin{IEEEbiographynophoto}{Alamgir Kabir}
\edit{PHOTO REQUIRED: kabir.jpg, 1$\times$1.25 in, 300 dpi}
\edit{received the Ph.D. degree in physics from \ldots{}, in YYYY. He is
currently \ldots{} with the Department of Physics, University of Dhaka, Dhaka,
Bangladesh. His research interests include \ldots{}.}
\end{IEEEbiographynophoto}

\end{document}
